\documentclass[a4paper,10pt]{report}\input{../0.Préambule/Préambule_cours_prof}\begin{document}

\newpage
\section*{Transformations planes - Agrandissement et réduction}
\addcontentsline{toc}{section}{Transformations planes - Agrandissement et réduction}

\subsection*{Utiliser une réduction ou un agrandissement}
\addcontentsline{toc}{subsection}{Utiliser une réduction ou un agrandissement}

\begin{exer1}

\vspace{8mm}

\includegraphics[scale=1]{../40.Image/15. Exo-Transformations planes - Agrandissement et réduction 1.png}

\vspace{-4.4cm}
\begin{flushright}
\includegraphics[scale=1]{../40.Image/15. Exo-Transformations planes - Agrandissement et réduction 2.png}
\end{flushright}

\vspace{-4cm}
\hspace{2cm} Indique si chaque silhouette correspond à une réduction, \\à un agrandissement ou à une déformation de la silhouette de chat ci-dessous.
\end{exer1}

\vspace{2.5cm}
\begin{exer1}
Dans chaque cas, la figure $2$ est-elle un agrandissement ou une réduction de la figure $1$ ? Justifie ta réponse

\includegraphics[scale=0.7]{../40.Image/15. Exo-Transformations planes - Agrandissement et réduction 3.png}
\end{exer1}

\noindent \begin{minipage}{9cm}
\begin{exer1}
Les figures $2$ et $3$ sont un agrandissement et \\une réduction de la figure $1$. Termine-les.
\end{exer1}

\begin{exer1}
Soit un triangle $ABC$ tel que $\widehat{ABC}=70$\degre ; \\$\widehat{BAC}=53$\degre et $AB =14$m. \\Construis une réduction de rapport $\dfrac{1}{200}$ de ce triangle.
\end{exer1}

\begin{exer1}
L'aire de la base d'un cylindre est de $51$cm$^2$. Quelle est l'aire de la base du cylindre obtenu après réduction de rapport $0,6$ ? Quel est son rayon, au dixième près ?
\end{exer1}

\begin{exer1}
Une figure a une aire de $124$cm$^2$. Après une réduction, on obtient une nouvelle figure dont l'aire est $89,59$cm$^2$. Détermine le rapport de réduction.
\end{exer1}
\end{minipage}
\hspace{5mm}
\begin{minipage}{7cm}
\begin{flushright}
\includegraphics[scale=0.8]{../40.Image/15. Exo-Transformations planes - Agrandissement et réduction 4.png}
\end{flushright}
\end{minipage}

\begin{exer1}
Un triangle $A'B'C'$ rectangle en $A'$ et d'aire $27$cm$^2$ est un agrandissement d'un triangle $ABC$, rectangle en $A$ tel que $AB=3$cm et $AC=2$cm. Calcule les longueurs $A'B'$ et $A'C'$.
\end{exer1}

\begin{exer2}
La pyramide du Louvre est un pyramide régulière à base carrée de coté $35$m  et de hauteur $22$m.
\begin{enumerate}[font=\color{exer2}]
\item Fais un schéma
\item Calcule le volume $\mathcal{V}$ de cette pyramide. Donne la valeur exacte en m$^3$ puis la valeur arrondie à l'unité.
\item Sur une maquette, on construit une réduction de cette pyramide, le coté de la base carrée mesure $7$cm. Calcule le coefficient de réduction.
\item Déduis-en le volume $\mathcal{V'}$ de la pyramide sur la maquette. donne la valeur exacte en cm$^3$ puis la valeur arrondie à l'unité.
\end{enumerate}
\end{exer2}

\begin{exer2}
On coupe une pyramide à mi-hauteur par un plan parallèle à la base.
\begin{enumerate}[font=\color{exer2}]
\item Exprime le volume $\mathcal{V'}$ de la petite pyramide en fonction du volume $\mathcal{V}$ de la pyramide de départ.
\item Montre que le volume $\mathcal{V''}$ du tronc de pyramide obtenu est égal aux $\dfrac{7}{8}$ du volume $\mathcal{V}$ de la pyramide de départ.
\end{enumerate}
\end{exer2}

\begin{exer2}
Une petite sphère a pour rayon $r$. Une grande sphère a pour rayon $R=3r$. Soient $v$ le volume de la petite sphère et $\mathcal{V}$ le volume de la grand sphère. Exprime $\mathcal{V}$ en fonction de $v$.
\end{exer2}
\end{document}